\documentclass[10pt,a4paper]{amsart}
\usepackage{amsmath,amssymb,amsthm,graphicx}
% Retained preamble for the Stacks Project chapter snapshot in this directory.
% The chapter's own preamble.tex is a copy of the shared Stacks Project preamble; the block
% below is sliced MECHANICALLY, line for line, from it. Nothing was added, removed or
% redefined; the wrapper class, the xr-hyper/hyperref wiring and the externaldocument list
% are omitted because the assembler supplies the document class.
%
% lines taken from preamble.tex: 3, 5, 9, 11, 12, 156, 157, 158, ... (51 lines total)
%
% package coverage restored from the chapter's own preamble (2026-10-13): xy, xy-2cell

\usepackage[all]{xy}
\xyoption{2cell}
\UseAllTwocells

\newenvironment{reference}{\comment}{\endcomment}
\newenvironment{slogan}{\comment}{\endcomment}
\newenvironment{history}{\comment}{\endcomment}
\theoremstyle{plain}
\newtheorem{theorem}[subsection]{Theorem}
\newtheorem{proposition}[subsection]{Proposition}
\newtheorem{lemma}[subsection]{Lemma}
\theoremstyle{definition}
\newtheorem{definition}[subsection]{Definition}
\newtheorem{example}[subsection]{Example}
\newtheorem{exercise}[subsection]{Exercise}
\newtheorem{situation}[subsection]{Situation}
\theoremstyle{remark}
\newtheorem{remark}[subsection]{Remark}
\newtheorem{remarks}[subsection]{Remarks}
\numberwithin{equation}{subsection}
\def\lim{\mathop{\mathrm{lim}}\nolimits}
\def\colim{\mathop{\mathrm{colim}}\nolimits}
\def\Spec{\mathop{\mathrm{Spec}}}
\def\Hom{\mathop{\mathrm{Hom}}\nolimits}
\def\Ext{\mathop{\mathrm{Ext}}\nolimits}
\def\SheafHom{\mathop{\mathcal{H}\!\mathit{om}}\nolimits}
\def\SheafExt{\mathop{\mathcal{E}\!\mathit{xt}}\nolimits}
\def\Sch{\mathit{Sch}}
\def\Mor{\mathop{\mathrm{Mor}}\nolimits}
\def\Ob{\mathop{\mathrm{Ob}}\nolimits}
\def\Sh{\mathop{\mathit{Sh}}\nolimits}
\def\NL{\mathop{N\!L}\nolimits}
\def\CH{\mathop{\mathrm{CH}}\nolimits}
\def\proetale{{pro\text{-}\acute{e}tale}}
\def\etale{{\acute{e}tale}}
\def\QCoh{\mathit{QCoh}}
\def\Ker{\mathop{\mathrm{Ker}}}
\def\Im{\mathop{\mathrm{Im}}}
\def\Coker{\mathop{\mathrm{Coker}}}
\def\Coim{\mathop{\mathrm{Coim}}}
\DeclareMathSymbol{\boxtimes}{\mathbin}{AMSa}{"02}
\def\QCohstack{\mathcal{QC}\!\mathit{oh}}
\def\Cohstack{\mathcal{C}\!\mathit{oh}}
\def\Spacesstack{\mathcal{S}\!\mathit{paces}}
\def\Quotfunctor{\mathrm{Quot}}
\def\Hilbfunctor{\mathrm{Hilb}}
\def\Curvesstack{\mathcal{C}\!\mathit{urves}}
\def\Polarizedstack{\mathcal{P}\!\mathit{olarized}}
\def\Complexesstack{\mathcal{C}\!\mathit{omplexes}}
\def\Pic{\mathop{\mathrm{Pic}}\nolimits}
\def\Picardstack{\mathcal{P}\!\mathit{ic}}
\def\Picardfunctor{\mathrm{Pic}}
\def\Deformationcategory{\mathcal{D}\!\mathit{ef}}
\usepackage{cleveref}
\crefname{originallemma}{original lemma}{original lemmas}
\Crefname{originallemma}{Original lemma}{Original lemmas}
\makeatletter
\newlabel{lemma-pic-gerbe-over-pic-functor}{{0DME}{source}{}{sourcepointer-1}{}}
\newlabel{lemma-pic-gerbe-over-pic-functor@cref}{{[originallemma][0][]0DME}{[1][0][]source}{}{}{}}
\makeatother
\crefname{originallemma}{original lemma}{original lemmas}
\Crefname{originallemma}{Original lemma}{Original lemmas}
\makeatletter
\newlabel{more-algebra-lemma-lift-pseudo-coherent-from-residue-field}{{0BCC}{source}{}{sourcepointer-2}{}}
\newlabel{more-algebra-lemma-lift-pseudo-coherent-from-residue-field@cref}{{[originallemma][0][]0BCC}{[1][0][]source}{}{}{}}
\makeatother
\crefname{originallemma}{original lemma}{original lemmas}
\Crefname{originallemma}{Original lemma}{Original lemmas}
\makeatletter
\newlabel{quot-lemma-flat-geometrically-connected-fibres}{{0D27}{source}{}{sourcepointer-3}{}}
\newlabel{quot-lemma-flat-geometrically-connected-fibres@cref}{{[originallemma][0][]0D27}{[1][0][]source}{}{}{}}
\makeatother
\crefname{originallemma}{original lemma}{original lemmas}
\Crefname{originallemma}{Original lemma}{Original lemmas}
\makeatletter
\newlabel{spaces-morphisms-lemma-valuative-criterion-separatedness}{{03KV}{source}{}{sourcepointer-4}{}}
\newlabel{spaces-morphisms-lemma-valuative-criterion-separatedness@cref}{{[originallemma][0][]03KV}{[1][0][]source}{}{}{}}
\makeatother
\crefname{originallemma}{original lemma}{original lemmas}
\Crefname{originallemma}{Original lemma}{Original lemmas}
\makeatletter
\newlabel{spaces-over-fields-lemma-characterize-trivial-pic-integral}{{0DMZ}{source}{}{sourcepointer-5}{}}
\newlabel{spaces-over-fields-lemma-characterize-trivial-pic-integral@cref}{{[originallemma][0][]0DMZ}{[1][0][]source}{}{}{}}
\makeatother
\crefname{originallemma}{original lemma}{original lemmas}
\Crefname{originallemma}{Original lemma}{Original lemmas}
\makeatletter
\newlabel{spaces-over-fields-lemma-proper-geometrically-reduced-global-sections}{{0DMY}{source}{}{sourcepointer-6}{}}
\newlabel{spaces-over-fields-lemma-proper-geometrically-reduced-global-sections@cref}{{[originallemma][0][]0DMY}{[1][0][]source}{}{}{}}
\makeatother
\crefname{originallemma}{original lemma}{original lemmas}
\Crefname{originallemma}{Original lemma}{Original lemmas}
\makeatletter
\newlabel{spaces-perfect-lemma-flat-proper-perfect-direct-image-general}{{0CTM}{source}{}{sourcepointer-7}{}}
\newlabel{spaces-perfect-lemma-flat-proper-perfect-direct-image-general@cref}{{[originallemma][0][]0CTM}{[1][0][]source}{}{}{}}
\makeatother
\crefname{originallemma}{original lemma}{original lemmas}
\Crefname{originallemma}{Original lemma}{Original lemmas}
\makeatletter
\newlabel{spaces-perfect-lemma-jump-loci}{{0D1Y}{source}{}{sourcepointer-8}{}}
\newlabel{spaces-perfect-lemma-jump-loci@cref}{{[originallemma][0][]0D1Y}{[1][0][]source}{}{}{}}
\makeatother
\crefname{originallemma}{original lemma}{original lemmas}
\Crefname{originallemma}{Original lemma}{Original lemmas}
\makeatletter
\newlabel{stacks-morphisms-lemma-gerbe-fppf}{{06QI}{source}{}{sourcepointer-9}{}}
\newlabel{stacks-morphisms-lemma-gerbe-fppf@cref}{{[originallemma][0][]06QI}{[1][0][]source}{}{}{}}
\makeatother
\crefname{originallemma}{original lemma}{original lemmas}
\Crefname{originallemma}{Original lemma}{Original lemmas}
\makeatletter
\newlabel{stacks-morphisms-lemma-lift-valuation-ring-through-flat-morphism}{{0DN5}{source}{}{sourcepointer-10}{}}
\newlabel{stacks-morphisms-lemma-lift-valuation-ring-through-flat-morphism@cref}{{[originallemma][0][]0DN5}{[1][0][]source}{}{}{}}
\makeatother
\title{The Picard functor over an integral-fibre family is separated}
\author{The Stacks Project Authors}
\date{stacks-project:moduli.tex:local-snapshot-20261012; GNU FDL 1.2 or later}
\begin{document}
\maketitle
\noindent\textit{Editorial note.} Statement span is lines 1713-1718 of sources/stacks-moduli/pinned-original-source.tex; proof span is lines 1720-1788, the complete author proof (67 body lines).\par\medskip
\noindent\textit{Editorial note.} Navigation stubs: 10 cited-but-not-included labels are supplied as non-load-bearing pointers with their pinned tags where the snapshot tag index has them.\par\medskip
\noindent\textit{Editorial note.} extra\_packages is empty on purpose: the retained preamble already loads xy with options and a second declaration clashes in the assembler wrapper.\par\medskip
\noindent\textit{Editorial note.} Text is retained verbatim, including the author's own typos and placeholders.\par\medskip
\noindent\textit{Editorial note.} Queue index 2 of queue\_moduli.json (pb=3038, stmt 1713-1718, proof 1720-1788).\par\medskip
\noindent\textit{Editorial note.} stacks-moduli wave-1: fresh adjudication accept (2026-10-14); pool per Chief order 08:5x (dpa superseded to lane-s9); moduli.tex 53 labels/48 with proof judged 25 accept / 23 reject / 5 pb=0 env-dup rows excluded; spans cite-free and reference/slogan/history-free so no patchenv needed.\par\medskip
\noindent\textit{Allowed standard prerequisites (editorial).}\begin{itemize}
\item Explicit prerequisite boundary. The cited inputs are prerequisites; the argument in the retained span is not replaced by any cited result.
\end{itemize}\medskip
\ifdefined\hypertarget\hypertarget{sourcepointer-1}{}\fi
\noindent\textit{Source reference.} Original lemma 0DME: sources/stacks-moduli/pinned-original-source.tex, source lines 1639--1642. Cited by the retained statement or proof but not part of the collected spans; supplied only as a navigation pointer. The argument does not depend on its content.\par\medskip
\ifdefined\hypertarget\hypertarget{sourcepointer-2}{}\fi
\noindent\textit{Source reference.} Original lemma 0BCC: sources/stacks-moduli/pinned-original-source.tex, source lines 1788--1788. Cited by the retained statement or proof but not part of the collected spans; supplied only as a navigation pointer. The argument does not depend on its content.\par\medskip
\ifdefined\hypertarget\hypertarget{sourcepointer-3}{}\fi
\noindent\textit{Source reference.} Original lemma 0D27: sources/stacks-moduli/pinned-original-source.tex, source lines 1788--1788. Cited by the retained statement or proof but not part of the collected spans; supplied only as a navigation pointer. The argument does not depend on its content.\par\medskip
\ifdefined\hypertarget\hypertarget{sourcepointer-4}{}\fi
\noindent\textit{Source reference.} Original lemma 03KV: sources/stacks-moduli/pinned-original-source.tex, source lines 1788--1788. Cited by the retained statement or proof but not part of the collected spans; supplied only as a navigation pointer. The argument does not depend on its content.\par\medskip
\ifdefined\hypertarget\hypertarget{sourcepointer-5}{}\fi
\noindent\textit{Source reference.} Original lemma 0DMZ: sources/stacks-moduli/pinned-original-source.tex, source lines 1788--1788. Cited by the retained statement or proof but not part of the collected spans; supplied only as a navigation pointer. The argument does not depend on its content.\par\medskip
\ifdefined\hypertarget\hypertarget{sourcepointer-6}{}\fi
\noindent\textit{Source reference.} Original lemma 0DMY: sources/stacks-moduli/pinned-original-source.tex, source lines 1788--1788. Cited by the retained statement or proof but not part of the collected spans; supplied only as a navigation pointer. The argument does not depend on its content.\par\medskip
\ifdefined\hypertarget\hypertarget{sourcepointer-7}{}\fi
\noindent\textit{Source reference.} Original lemma 0CTM: sources/stacks-moduli/pinned-original-source.tex, source lines 1788--1788. Cited by the retained statement or proof but not part of the collected spans; supplied only as a navigation pointer. The argument does not depend on its content.\par\medskip
\ifdefined\hypertarget\hypertarget{sourcepointer-8}{}\fi
\noindent\textit{Source reference.} Original lemma 0D1Y: sources/stacks-moduli/pinned-original-source.tex, source lines 1788--1788. Cited by the retained statement or proof but not part of the collected spans; supplied only as a navigation pointer. The argument does not depend on its content.\par\medskip
\ifdefined\hypertarget\hypertarget{sourcepointer-9}{}\fi
\noindent\textit{Source reference.} Original lemma 06QI: sources/stacks-moduli/pinned-original-source.tex, source lines 1788--1788. Cited by the retained statement or proof but not part of the collected spans; supplied only as a navigation pointer. The argument does not depend on its content.\par\medskip
\ifdefined\hypertarget\hypertarget{sourcepointer-10}{}\fi
\noindent\textit{Source reference.} Original lemma 0DN5: sources/stacks-moduli/pinned-original-source.tex, source lines 1788--1788. Cited by the retained statement or proof but not part of the collected spans; supplied only as a navigation pointer. The argument does not depend on its content.\par\medskip
\begin{lemma}
\label{lemma-pic-functor-uniqueness-part}
Assume the geometric fibres of $X \to B$ are integral
in addition to the other assumptions in this section.
Then $\Picardfunctor_{X/B} \to B$ is separated.
\end{lemma}

\begin{proof}
Since $\Picardfunctor_{X/B} \to B$ is quasi-separated, it suffices
to check the uniqueness part of the valuative criterion, see
Morphisms of Spaces, Lemma
\ref{spaces-morphisms-lemma-valuative-criterion-separatedness}.
This immediately reduces to the following problem: given
\begin{enumerate}
\item a valuation ring $R$ with fraction field $K$,
\item an algebraic space $X$ proper and flat over $R$
with integral geometric fibre,
\item an element $a \in \Picardfunctor_{X/R}(R)$ with
$a|_{\Spec(K)} = 0$,
\end{enumerate}
then we have to prove $a = 0$. Applying
Morphisms of Stacks, Lemma
\ref{stacks-morphisms-lemma-lift-valuation-ring-through-flat-morphism}
to the surjective flat morphism
$\Picardstack_{X/R} \to \Picardfunctor_{X/R}$
(surjective and flat by Lemma \ref{lemma-pic-gerbe-over-pic-functor} and
Morphisms of Stacks, Lemma \ref{stacks-morphisms-lemma-gerbe-fppf})
after replacing $R$ by an extension we may assume
$a$ is given by an invertible $\mathcal{O}_X$-module
$\mathcal{L}$. Since $a|_{\Spec(K)} = 0$ we find
$\mathcal{L}_K \cong \mathcal{O}_{X_K}$ by
Quot, Lemma \ref{quot-lemma-flat-geometrically-connected-fibres}.

\medskip\noindent
Denote $f : X \to \Spec(R)$ the structure morphism.
Let $\eta, 0 \in \Spec(R)$ be the generic and closed point.
Consider the perfect complexes
$K = Rf_*\mathcal{L}$ and $M = Rf_*(\mathcal{L}^{\otimes -1})$
on $\Spec(R)$, see Derived Categories of Spaces, Lemma
\ref{spaces-perfect-lemma-flat-proper-perfect-direct-image-general}.
Consider the functions
$\beta_{K, i}, \beta_{M, i} : \Spec(R) \to \mathbf{Z}$
of Derived Categories of Spaces, Lemma
\ref{spaces-perfect-lemma-jump-loci} associated to $K$ and $M$.
Since the formation of $K$ amd $M$ commutes with
base change (see lemma cited above) we find
$\beta_{K, 0}(\eta) = \beta_{M, 0}(\eta) = 1$ by
Spaces over Fields, Lemma
\ref{spaces-over-fields-lemma-proper-geometrically-reduced-global-sections}
and our assumption on the fibres of $f$.
By upper semi-continuity we find
$\beta_{K, 0}(0) \geq 1$ and $\beta_{M, 0}(0) \geq 1$.
By 
Spaces over Fields, Lemma
\ref{spaces-over-fields-lemma-characterize-trivial-pic-integral}
we conclude that the restriction of $\mathcal{L}$
to the special fibre $X_0$ is trivial. In turn this gives
$\beta_{K, 0}(0) = \beta_{M, 0} = 1$ as above.
Then by More on Algebra, Lemma
\ref{more-algebra-lemma-lift-pseudo-coherent-from-residue-field}
we can represent $K$ by a complex of the form
$$
\ldots \to 0 \to R \to R^{\oplus \beta_{K, 1}(0)} \to
R^{\oplus \beta_{K, 2}(0)} \to \ldots
$$
Now $R \to R^{\oplus \beta_{K, 1}(0)}$ is zero
because $\beta_{K, 0}(\eta) = 1$. In other words
$K = R \oplus \tau_{\geq 1}(K)$ in $D(R)$ where $\tau_{\geq 1}(K)$
has tor amplitude in $[1, b]$ for some $b \in \mathbf{Z}$.
Hence there is a global section $s \in H^0(X, \mathcal{L})$
whose restriction $s_0$
to $X_0$ is nonvanishing (again because formation of $K$
commutes with base change). Then $s : \mathcal{O}_X \to \mathcal{L}$
is a map of invertible sheaves whose restriction to $X_0$
is an isomorphism and hence is an isomorphism as desired.
\end{proof}

\end{document}
